% answers_short.tex —— 2021 年·新高考Ⅱ卷 · 简短答案速查
% 由 tools/make_short_answers.py 自 answers.tex 的 \gksolhead/\gkanswer 提取，
% 并与 2021/2021_数学真题_新高考Ⅱ_答案要点.md 逐题比对后生成，勿手改（改源 MD 后重跑）。

\begin{gknotice}
\textbf{使用说明}：本节只给最终答案便于快速核对，完整推导见解析版。答案为编写者独立推导，\textbf{非官方标准答案}。
\end{gknotice}

\gkpart{单选题}
\begin{center}\begin{tabular}{ccccccccc}
\toprule
题号 & 1 & 2 & 3 & 4 & 5 & 6 & 7 & 8 \\
\midrule
答案 & A & B & B & C & D & D & C & B \\
\bottomrule
\end{tabular}\end{center}

\gkpart{多选题}
\begin{center}\begin{tabular}{ccccc}
\toprule
题号 & 9 & 10 & 11 & 12 \\
\midrule
答案 & AC & BC & ABD & ACD \\
\bottomrule
\end{tabular}\end{center}

\gkpart{填空题}
\begin{itemize}[leftmargin=2.6em]
  \item[\textbf{13.}] $y=\pm\sqrt3\,x$
  \item[\textbf{14.}] $f(x)=x^{2}$（答案不唯一，一切 $x^{2k},k\in\mathbf N^{*}$ 均可）
  \item[\textbf{15.}] $-\dfrac92$
  \item[\textbf{16.}] $(0,1)$
\end{itemize}

\gkpart{解答题}
\begin{itemize}[leftmargin=2.6em]
  \item[\textbf{17.}] (1) $a_n=2n-6$；(2) $n_{\min}=7$
  \item[\textbf{18.}] (1) $S=\dfrac{15\sqrt7}4$；(2) 存在，$a=2$
  \item[\textbf{19.}] (1) 见解析版；(2) $\cos\theta=\dfrac23$
  \item[\textbf{20.}] (1) $\dfrac{x^{2}}3+y^{2}=1$；(2) 充要条件成立（两向都要证）
  \item[\textbf{21.}] (1) $E(X)=1$；(2) 见解析版；(3) 见解析版
  \item[\textbf{22.}] (1) 见解析版；(2) 两个条件下 $f(x)$ 均恰有一个零点
\end{itemize}
